\documentclass[preprint, 3p,
authoryear]{elsarticle} %review=doublespace preprint=single 5p=2 column
%%% Begin My package additions %%%%%%%%%%%%%%%%%%%

\usepackage[hyphens]{url}

  \journal{An awesome journal} % Sets Journal name

\usepackage{graphicx}
%%%%%%%%%%%%%%%% end my additions to header

\usepackage[T1]{fontenc}
\usepackage{lmodern}
\usepackage{amssymb,amsmath}
% TODO: Currently lineno needs to be loaded after amsmath because of conflict
% https://github.com/latex-lineno/lineno/issues/5
\usepackage{lineno} % add
\usepackage{ifxetex,ifluatex}
\usepackage{fixltx2e} % provides \textsubscript
% use upquote if available, for straight quotes in verbatim environments
\IfFileExists{upquote.sty}{\usepackage{upquote}}{}
\ifnum 0\ifxetex 1\fi\ifluatex 1\fi=0 % if pdftex
  \usepackage[utf8]{inputenc}
\else % if luatex or xelatex
  \usepackage{fontspec}
  \ifxetex
    \usepackage{xltxtra,xunicode}
  \fi
  \defaultfontfeatures{Mapping=tex-text,Scale=MatchLowercase}
  \newcommand{\euro}{€}
\fi
% use microtype if available
\IfFileExists{microtype.sty}{\usepackage{microtype}}{}
\usepackage[]{natbib}
\bibliographystyle{elsarticle-harv}

\ifxetex
  \usepackage[setpagesize=false, % page size defined by xetex
              unicode=false, % unicode breaks when used with xetex
              xetex]{hyperref}
\else
  \usepackage[unicode=true]{hyperref}
\fi
\hypersetup{breaklinks=true,
            bookmarks=true,
            pdfauthor={},
            pdftitle={Linking calculated and perceved accessibility to suggest equity standards},
            colorlinks=false,
            urlcolor=blue,
            linkcolor=magenta,
            pdfborder={0 0 0}}

\setcounter{secnumdepth}{5}
% Pandoc toggle for numbering sections (defaults to be off)


% tightlist command for lists without linebreak
\providecommand{\tightlist}{%
  \setlength{\itemsep}{0pt}\setlength{\parskip}{0pt}}




\usepackage{pdflscape}
\usepackage{booktabs}
\usepackage{longtable}
\usepackage{array}
\usepackage{multirow}
\usepackage{wrapfig}
\usepackage{float}
\usepackage{colortbl}
\usepackage{pdflscape}
\usepackage{tabu}
\usepackage{threeparttable}
\usepackage{threeparttablex}
\usepackage[normalem]{ulem}
\usepackage{makecell}
\usepackage{xcolor}



\begin{document}


\begin{frontmatter}

  \title{Linking calculated and perceved accessibility to suggest equity
standards}
    \author[Some Institute of Technology]{Alice Anonymous%
  \corref{cor1}%
  \fnref{1}}
   \ead{alice@example.com} 
    \author[Another University]{Bob Security%
  %
  }
   \ead{bob@example.com} 
    \author[Another University]{Cat Memes%
  %
  \fnref{2}}
   \ead{cat@example.com} 
    \author[Some Institute of Technology]{Derek Zoolander%
  %
  \fnref{2}}
   \ead{derek@example.com} 
      \affiliation[Some Institute of Technology]{
    organization={Big Wig University},addressline={1 main
street},city={Gotham},postcode={123456},state={State},country={United
States},}
    \affiliation[Another University]{
    organization={Department},addressline={A street
29},city={Manchester,},postcode={2054 NX},country={The Netherlands},}
    \cortext[cor1]{Corresponding author}
    \fntext[1]{This is the first author footnote.}
    \fntext[2]{Another author footnote.}
  
  \begin{abstract}
  This is the abstract.

  It consists of two paragraphs.
  \end{abstract}
    \begin{keyword}
    keyword1 \sep 
    keyword2
  \end{keyword}
  
 \end{frontmatter}

\section{Introduction}\label{introduction}

Transportation equity refers to the distribution of the burdens and
benefits of the transportation system among the population (Karner et
al., 2020, 2024; Pereira et al., 2017; Pereira \& Karner, 2021). Burdens
include health-related concerns arising from exposure to air pollution,
noise, and traffic; sedentary lifestyles; and the risk of injury or
death from collisions. In contrast, the primary benefit is the
improvement of population accessibility (Pereira \& Karner, 2021),
understood as the ease of reaching (spatially distributed) opportunities
and destinations (5--8), such as employment, healthcare, educational
services, and other places for social interaction.

Accessibility is one of the primary outcomes of spatial development
(Páez et al., 2012). It is influenced by the spatial distribution of the
population and opportunities and depends on the characteristics and
performance of the transportation network, as a product of both
transportation systems and land use patterns (Wee \& Geurs, 2011). There
is growing consensus among researchers and transportation planners that
the goal of transport policy should be to improve accessibility. This
contrasts with traditional planning approaches that have prioritized
mobility patterns that stimulate private automobile use and higher
speeds as solutions to congestion and long commuting times.
Mobility-centered policies often focus on quantitative metrics, such as
increasing the number of trips, raising average speeds, and reducing
congestion, treating mobility as an end in itself. However, this
approach is problematic because people rarely travel simply for the sake
of moving; instead, they travel to access opportunities, activities, or
other people at their destinations. Therefore, since the main goal is to
enable people to access activities, land use and transport planning
practices must be redesigned to prioritize accessibility.

Insufficient or inadequate accessibility can prevent individuals from
participating in location-based activities they wish to engage in for
reasons beyond their control, in a process described in the literature
as Transport-Related Social Exclusion (TRSE) (Casas, 2007; Lucas, 2012;
Paez et al., 2009; Preston \& Rajé, 2007; Luz \& Portugal, 2022).

Insufficient accessibility, when combined with other forms of social and
economic disadvantage, can result in transportation poverty (Casas,
2007; Litman, 2024; Lucas, 2012, 2019; Paez et al., 2009; Preston \&
Rajé, 2007): the failure of transportation and land use systems to
provide access to essential services or employment (Luz \& Portugal,
2022). Transportation poverty increases vulnerability among individuals
who already face challenges associated with age (van Dülmen et al.,
2022), income (Allen \& Farber, 2019), race (Awaworyi Churchill, 2020),
gender (Murphy et al., 2023), disability (Upham et al., 2022), among
other factors (Kim et al., 2024).

With the gentrification of Canadian downtowns and the suburbanization of
low-income populations, it is estimated that over one million Canadians
live in transport poverty (Allen \& Farber, 2019). An equity standard,
when operationalized, defines the level of fairness to be adopted and
the acceptable inequality (Soukhov, Anastasia et al., 2024). However,
the lack of standards for the equitable distribution of transportation
benefits complicates efforts to address transportation poverty.

Additionally, there is an absence of transportation data in the form of
open data products to support urban planners (Arribas-Bel et al., 2021;
Behrendt \& Sheller, 2021; Páez, 2021), and transportation planners
receive limited formal guidance in assessing the equity implications of
their projects (Miller, 2018; Palm et al., 2023; te Brömmelstroet et
al., 2016). Lastly, few studies explicitly link accessibility levels
with employment outcomes, such as employment opportunities declined or
accepted due to transport conditions (Soukhov, Anastasia et al., 2024),
or individuals' perceptions - for instance, their opinion on how easy is
to reach employment opportunities and whether jobs are available within
an ideal commuting time. Combined, this results in a gap in the
development and implementation of accessibility standards to ensure
equity in transportation planning and policy.

This study is guided by the following research questions: How can
calculated job accessibility be related to perceived accessibility and
employment outcomes? And could a given level of accessibility serve as a
standard or goal for equitable transport planning? To answer these
questions, this study has the general objective of identifying key
thresholds in the relationship between calculated job accessibility and
people's perceptions of employment outcomes in order to define equity
standards for transportation planning.

Our research advances the current understanding of the literature on the
equitable distribution of transportation benefits by proposing a method
that identifies calculated accessibility thresholds based on their
relationship with perceived accessibility. Identifying minimum job
accessibility levels in practice would facilitate the implementation of
normative criteria that can guarantee standards to prevent
equity-deserving groups from being excluded from the labour market.
Furthermore, accessibility thresholds can provide policymakers with an
understanding of how much accessibility is needed to meet the
fundamental needs of the population, increasing the potential for
government adoption and enhancing the comprehensibility and relevance of
standards for community members and those experiencing unequal mobility.

We also present, in a transparent manner, the moral guidelines that
support the distribution of accessibility. In addition, to facilitate
collaboration and further analysis and to ensure the transparency and
replicability of our study, we developed this paper using literate
programming, with the data analysis code accessible through our GitHub
page (\url{https://github.com/dias-bruno/NSTP}), in alignment with best
practices in spatial data science (Arribas-Bel et al., 2021).

\section{Theoretical Background}\label{theoretical-background}

\subsection{Calculated and perceived
accessibility}\label{calculated-and-perceived-accessibility}

The most commonly used accessibility measures are based on the gravity
model of Hansen and make use of spatial data. This model considers the
number of opportunities that a person can reach to be proportional to
the sum of opportunities in each destination, discounted by the travel
cost to reach that destination. This type of metric is classified as a
location-based measure, representing a characteristic of a given
location that assumes that all people in this location have equal access
to the spatially distributed opportunities. Such measures are sensitive
to land use factors and to the performance of the transportation system
(accessibility 2004). However, they do not take into account the
individual characteristics of people.

The perceived accessibility, on the other hand, refers to the perception
of people of how easily they can reach opportunities (Lattman et al.,
2016; Pot et al., 2021), resulting from not only the interaction of
land-use and transport infrastructure, but also from travel options
(e.g., having a car or a transit pass), travel experiences, and
individual characteristics and preferences (e.g., Lattman et al., 2016;
Mehdizadeh 2025). Often, perceived accessibility deviates from
place-based calculated accessibility, with some research indicating that
calculated accessibility might not be able to capture the multiple ways
in which people experience a lack of access to opportunities (Zhu \&
Lucas, 2025; Paromita, 2026). Explanations for these differences are
partially attributed to differences in individuals' travel attitudes,
capabilities, and preferences (J. De Vos et al.~2025 -
\url{https://www-sciencedirect-com.libaccess.lib.mcmaster.ca/science/article/pii/S0965856425001818?via\%3Dihub};
van der Vlugt et al., 2022). .

Even though, as demonstrated in the following subsection, calculated
accessibility is still commonly used to evaluate the effect that the
interaction between land-use characteristics and the transportation
system has on employment outcomes. In addition, existing studies on
perceived accessibility seem to have difficulties in creating
standardized definitions and measurement tools, contributing to the
somewhat limited knowledge on the topic (e.g., Negm et al., 2025). In
this paper, as suggested by Pot et al.~(2024), we define sufficient
levels of accessibility based on the relationship between calculated
accessibility and people's perception of employment outcomes, while also
controlling for individual characteristics.

\subsection{Job accessibility and employment
outcomes}\label{job-accessibility-and-employment-outcomes}

Studies connecting employment outcomes to job accessibility date back to
the 1960s, when John Kain (1968) proposed the spatial mismatch
hypothesis while analyzing Detroit and Chicago, USA. The author
hypothesized that the low employment rates among Black Americans in
inner cities resulted from employment decentralization and the
difficulty faced by the Black population in moving out of these areas,
primarily due to racial discrimination and segregation.

Since the work of Kain , studies linking calculated job accessibility
and different employment outcomes have presented mixed results. Bania et
al.~(2008) found no significant effect for job access when examining its
relationship with employment, earnings, hourly wage, and weekly hours
worked for women leaving welfare in Ohio, USA. Delmelle et al.~(2021)
also found that changes in car accessibility had no significant impact
on employment rates for different population income groups in the
Metropolitan Area of Charlotte, North Carolina, USA, arguing that the
spatial separation between lower-wage workers and their workplaces might
already be relatively small. Studying Los Angeles, USA, Merlin and Hu
(2017) found mixed-effects of job accessibility on employment status,
with a stronger association for more highly educated groups. The authors
argued that lower-educated groups may be more limited by non-spatial
factors that impede their employment opportunities, such as
qualifications and social skills. On the other hand, for workers with a
higher level of education, the time and monetary costs of spatial
separation itself may play a relatively more important role, so
accessibility measures may have a greater impact on their decisions and
ability to enter the labor market.

In contrast, other studies found significant effects of job
accessibility on employment outcomes (Bastiaanssen et al., 2021; Duarte
et al., 2023; Jin \& Paulsen, 2018; Johnson et al., 2017; Matas et al.,
2010). Jin and Paulsen (2017) used a difference-in-differences model for
the Chicago Metropolitan Area, finding that increasing job accessibility
between 2010 and 2017 was associated with reduced unemployment rates.
Duarte et al.~(2023) found that job accessibility had a positive effect
on participation in the labour market for women, a negative effect on
informal employment for women, and a positive effect on the probability
of being employed for both women and men. More recently, Bastiaanssen et
al.~(2025) examined differential employment effects of job accessibility
by public transport in combination with the bicycle for different
educational groups at the national level in the Netherlands. They found
that jobs for lower-educated groups are often located outside areas of
high accessibility, thereby reducing job accessibility among these
groups, while jobs for higher-educated groups tend to be concentrated in
and around downtowns and are associated with greater accessibility.

\subsection{Equity Standards}\label{equity-standards}

Accessibility standards are often operationalized as a type of
opportunity-based standard, defined using thresholds in access to
opportunities (e.g., the number of opportunities within a 10-minute
walking trip) or based on travel impedance thresholds (e.g., a given
distance to a subway station). After a systematic review, Souhkov et
al.~identified that this type of standard is the most commonly applied
in transportation literature (68\%).

Yenisetty \& Bahadure (2020) set an opportunity-based standard,
considering a distance of more than 1,200 meters between amenities,
services, and facilities by public transport to be unfair when analyzing
cities in India. Tiznado-Aitken (2018) defined unfair cases as those
where bus stops and metro stations are located more than 1 kilometer
away. Sun \& Thakuriah adopted 400 meters as the maximum acceptable
walking distance to bus and tram stops, and 800 meters to ferry stations
in England. In another study, Peungnumsai et al.~(2020) identified
opportunity standards through the ratio of public transport supply to
demand, considering conditions fair when demand is greater than supply.

However, one must to consider the possible implications of relying on
normative values to define thresholds. For example, when adopting a
fixed distance to a certain bus stops or train stations the
justification for this choice is often unclear. In addition, the use of
distance-based thresholds may be problematic, since the time required to
cross a certain distance can vary considerably within a city, even for
the same mode of transportation.

\section{Materials and Methods}\label{materials-and-methods}

The methodology for obtaining standards of accessibility based on the
relationship between calculated measures and people's perceptions can be
synthesized in six main steps: modelling job accessibility, selecting
population groups with the lowest perceived accessibility, predicting
employment outcomes using calculated accessibility, setting equity
standards for transportation, and defining priority geographical areas
for policy interventions. We recognize accessibility as the transport
benefit to be distributed across society, but we also find it necessary
to define the moral principles guiding such distribution, how the
distribution should occur to ensure fairness, and the requirements for
considering a population group a priority group. Therefore, the
following guidelines have been defined:

\begin{enumerate}
\def\labelenumi{\arabic{enumi}.}
\tightlist
\item
  We recognize accessibility as a human capability, and a higher level
  of accessibility implies more choice.
\item
  A fair transportation and land-use system provides sufficient
  accessibility to all members of society. A minimum level of job
  accessibility is essential for individuals to meet their basic needs
  and, while not sufficient on its own, it is a prerequisite for
  promoting equality of job opportunities (sufficientarian perspective).
\item
  The minimum level of accessibility should be defined in terms of the
  most unsatisfied population group in relation to the transportation
  system, while considering equity perspectives and implications.
\item
  A fair distribution of accessibility should focus on reducing
  inequalities of opportunity by prioritizing disadvantaged populations,
  while simultaneously aiming to improve overall accessibility levels.
\item
  Although accessibility standards should be based on the needs of
  disadvantaged populations, it is also necessary to identify areas
  where public policy interventions should be targeted, as the needs of
  populations can vary significantly depending on individual
  constraints.
\end{enumerate}

\subsection{Study Area and Data
Sources}\label{study-area-and-data-sources}

Our study area refers to the urban conurbation formed by the Greater
Toronto Area and the Hamilton Area (GTHA), in Ontario, Canada. More
specifically, the study area refers to the Census Divisions of the
regional municipalities of Durham, Hamilton, Halton, Peel, Toronto, and
York. According to the latest census (2021), over 7.2 million residents
live in this area, which is spread across around 7,300 square
kilometers.

Regarding the primary data sources, this study uses data from the 2021
Census of Population (CoP), from Statistics Canada, and from the
National Survey of Transport Poverty (NSTP), from the Mobilizing Justice
Partnership. The Census was selected due to its various socio-economic
and demographic variables and due to the ``Commuting to Work''
questionnaire. The Census allows the identification of individual in the
labour force and include variables required to calculate job
accessibility, such as place of work, commuting modes, and commuting
duration. The Census also provides standardized spatial datasets with
statistical and administrative geographical boundaries, including the
Census Tract (CTs), which are used in this study as the spatial units
used for calculating accessibility measures. According to the Census,
3,868,120 individuals in the labour force reside in the study area,
accessing a total of 3,043,730 job opportunities. Figure 2 shows the
total labour force and the number of jobs across census tracts in the
study area.

The second survey, NSTP, is a large-scale survey designed to understand
transport-related social exclusion (TRSE) and transport poverty in
Canada, collecting responses from more than 27,000 people across the
country and providing an unprecedented look at the transportation
challenges Canadians face, especially among equity-deserving groups
whose experiences are often overlooked. The survey is openly available
to researchers, decision-makers, and the public, and offers information
on household dynamics, respondents' perceptions, and individual-level
socioeconomic data. The NSTP provides respondent weights to enable the
creation of statistics that are consistent and representative of the
Canadian population (Morency et al., 2024). This survey was selected
because it contains respondents' perceptions of accessibility and
employment outcomes, as well as socioeconomic and demographic variables
at the individual level, with a geographic resolution that enables the
linking with calculated accessibility.

Not all questions from the NSTP were answered by all survey respondents.
The questions were defined according to demographic and socioeconomic
profiles and to previous responses. For this study, the NSTP sample
corresponds to respondents residing in the GTHA who also answered the
questions displayed in Table \ref{tab:out-description}, totalling 2,371
respondents that represent 3,684,512 individuals in the population
(weighted sample). Using a continuum scale from 0 (completely disagree)
to 100 (completely agree), the respondents answered to what extent they
agreed with the statements in Table \ref{tab:out-description}.

\begingroup\fontsize{8}{10}\selectfont

\begin{longtable}[t]{lc}
\caption{\label{tab:unnamed-chunk-2}\label{tab:out-description}Weighted summary statistics of the outcome variables.}\\
\toprule
To what extent do you agree with the following statements? & \textbf{N = 3,684,512}\\
\midrule
There are too few jobs available for me in my ideal commuting time & 58.84 (27.71)\\
With the transport options available to me, I can easily reach work & 72.64 (22.58)\\
I have had to decline employment opportunities because of my transport situation & 41.06 (31.30)\\
It takes too long to travel to work & 54.04 (29.26)\\
\bottomrule
\multicolumn{2}{l}{\rule{0pt}{1em}\textsuperscript{1} Mean (SD)}\\
\end{longtable}
\endgroup{}

\subsection{Modelling accessibility}\label{modelling-accessibility}

In this study, we implement the multimodal spatial availability measure
to model job accessibility. This measure is an extension of the spatial
availability measure to include travel by different modes. Both measures
consider the competition for opportunities, which avoids a single job
opportunity being equally available to more than one person. Competitive
measures matter when considering employment outcomes because jobs are
exclusive, and their use by one person prevents their use by another. In
addition, these measures relies on a proportional allocation that
balances the cost of travel and the population size, imposing a
constrained calculation that allocates the opportunities to each origin
ensuring that, when all jobs allocated to an origin are summed, they
equal the total number of opportunities in the region. Lastly, both
measures have the advantage of resulting in an interpretable indicator,
in which the output of the model refers to the number of jobs reachable
in each spatial unit.

We calculated the multimodal spatial availability \(V_i^m\) using
Equation (\ref{eq:spatial-availability-multimodal}):

\begin{equation}
\label{eq:spatial-availability-multimodal}
V^m_{i} = \sum_{j=1}^J O_j\ F^{tm}_{ij}
\end{equation}

where:

\begin{itemize}
\tightlist
\item
  \(V^m_{i}\) is the spatial availability of opportunities at origin
  \(i\) and travel mode \(m\);
\item
  \(O_j\) is the number of opportunities at destination \(j\);
\item
  \(F^{tm}_{ij}\) is a balancing factor allocating opportunities at
  \(j\) to origin \(i\) considering mode \(m\).
\end{itemize}

\(F^{tm}_{ij}\) incorporates both population demand and travel
impedance, defined according to Equation
\ref{eq:multimodal-balancing-factors}:

\begin{equation}
\label{eq:multimodal-balancing-factors}
F^{tm}_{ij} = \frac{F^{pm}_{i} \cdot F^{cm}_{ij}}{\sum_{m=1}^M \sum_{i=1}^N F^{pm}_{i} \cdot F^{cm}_{ij}}
\end{equation}

The population-based allocation factor is obtained by:

\begin{itemize}
\tightlist
\item
  \(F^{pm}_{i} = \frac{P_{i}^m}{\sum_{m}\sum_{i} P_{i}^m}\)
\end{itemize}

where \(P_{i}^m\) represents the population at origin \(i\) using mode
\(m\).

On its turn, the travel impedance allocation factor is given by:

\begin{itemize}
\tightlist
\item
  \(F_{ij}^{cm} = \frac{f^m(c_{ij}^m)}{\sum_{m} \sum_{i} f^m(c_{ij}^m)}\)
\end{itemize}

where:

\begin{itemize}
\tightlist
\item
  \(c^m_{ij}\) is the cost to travel between \(i\) and \(j\) for mode
  \(m\);
\item
  \(f^m{c^m_{ij}}\) is an impedance function representing the effect of
  travel cost.
\end{itemize}

Using microdata from the 21CoP, we calculated for each CT the number of
people in the labour force that commute to work by walking, cycling,
car-motorized modes (those who commute by car, truck, or van as a driver
or passenger, and motorcycles or scooters), and public transit (those
who commute by bus, subway, light rail, streetcar, commuter train, and
passenger ferry). We distributed remote workers according to the share
of each transportation mode in each CT. Then, we summed the number of
jobs in each CT considering the place of work mentioned by each worker.

We calculated the travel times between origins and destinations for each
travel mode using the R package r5r with a street network obtained from
OpenStreetMap. We used the General Transit Feed Specification (GTFS)
obtained from the Canadian Public Transit Network Database to retrieve
the public transit schedule and coverage. Considering that our study
area contains different transit agencies, we evaluated the GTFS files to
identify a common date/week with public transit service. After this
analysis, we set 8 a.m. (morning peak time) of December 22, 2025, as the
trip departure time. We also used Digital Elevation Models (DEM) from
the Canadian Digital Surface Model (CDSM), with a spatial resolution of
about 20 meters, to estimate the effect of declivity on walking and
cycling times. We set a maximum duration of 180 minutes (3 hours) for
car and public transit trips, with a maximum walking time of 5 minutes
to reach the vehicle and 15 minutes to reach the public transit
stop/station; for walking and cycling trips, we set a maximum duration
of 90 minutes (one and a half hours). After estimating the travel times
between CTs, we assumed that all intra-CT trip durations are equal to
0.1 minute.

Commute duration was used to calibrate Probability Density Functions
(PDFs) and represent travel impedance, following the approach of Soukhov
and Páez \citeyearpar{soukhov2024}, Verma and Ukkusuri
\citeyearpar{sveda2023}, and Dos Santos et al.~{[}ACTIVE TRAVEL
PAPER{]}. We applied the fitdistrplus package
\citep{delignette2015fitdistrplus} and tested different functional forms
(uniform, negative exponential, gamma, normal, and lognormal
distributions) to select the best PDF for each mode of transportation
and municipality, using a selection criteria the lowest Akaike
Information Criterion (AIC) value \citep{akaike1974} (Appendix 1
displays the estimated impedance functions).

\subsection{Modelling employment
outcomes}\label{modelling-employment-outcomes}

We combined the multimodal spatial availability measures with variables
at the individual-level from the NTSP to model perceptions on employment
outcomes, assessing two types of weighted regression models. Considering
respondent weights, we transformed the continuum perceptions to a likert
scale based on their quintile position, categorizing them between:
``Most Negative'', ``Negative'', ``Neutral'', ``Positive'', and ``Most
Positive''. Weighted ordinal logistic regressions were modelled to
examine the effect of between the accessibility indicators on the
perception categories, using the Equation \ref{eq:ordinal-logit}:

\begin{equation}
\label{eq:ordinal-logit}
\log(\frac{P(Y\le j)}{P(Y\ge j)}) = \theta_j - \beta_1 X_1 + \beta_2 X_2 + ... + \beta_p X_p
\end{equation}

Where \(Y\) represents the perception category of the respondent
(ranging from ``Most Negative'' to ``Most Positive''), and
\(X_1, X2, ..., X_p\) represent the set of explanatory variables,
including multimodal spatial availability measures and individual and
household characteristic obtained from the NTSP. \(\theta_j\) and
\(\beta_p\) are unknown parameters to be estimated. The parameters
\(\beta_p\) are constant for all outcome categories and measure the
change in the log-oggs of reporting higher perception categories
associated with each explanatory variable, while \(\theta_j\) represents
the intercept separating the perception categories.

Table \ref{tab:var-description} presents the variables included in the
regression models and descriptive statistics (weighted mean and weighted
standard deviation for continuous and dummy variables; and weighted
count and percentage for categorical variables). We hypothesise that
higher values of calculated accessibility will increase the probability
of being in more positive perceptions, and that perceptions of
accessibility can be connected to the multimodal spatial availability
indicators. All other explanatory variables are employed as control
variables to properly measure the effect of accessibility indicators on
dependent variables.

\begingroup\fontsize{8}{10}\selectfont

\begin{longtable}[t]{lc}
\caption{\label{tab:vardescriptions}\label{tab:var-description}Weighted summary statistics of the explanatory variables.}\\
\toprule
Explanatory variables & \textbf{N = 3,684,512}\\
\midrule
Active transportation job accessibility per user & 0.06 (0.13)\\
Transit job accessibility per user & 0.14 (0.09)\\
Car job accessibility per user & 1.00 (0.27)\\
Active transportation user (dummy) & 0.15 (0.36)\\
Public transit user (dummy) & 0.23 (0.42)\\
\addlinespace
Number of Cars & 1.30 (0.92)\\
Number of Bikes & 1.15 (1.33)\\
Preferred mode match most frequent mode & 2,794,984 (76\%)\\
Indigenous or racialized person & 1,917,721 (52\%)\\
Temporary visa or refugee & 272,937 (7.4\%)\\
\addlinespace
Person with disability & 731,104 (20\%)\\
Woman, transgender or non-binary person & 1,828,742 (50\%)\\
Age & 40.37 (13.57)\\
Household size & 2.89 (1.61)\\
Household income & \\
\addlinespace
\hspace{1em}Less than \$40,000 & 583,922 (16\%)\\
\hspace{1em}\$40,000 to \$80,000 & 1,174,003 (32\%)\\
\hspace{1em}\$80,000 to \$120,000 & 965,106 (26\%)\\
\hspace{1em}More than \$120,000 & 961,480 (26\%)\\
Population density (people/km2) & 7,818.93 (13,582.70)\\
\addlinespace
Population centre class & \\
\hspace{1em}Rural & 92,135 (2.5\%)\\
\hspace{1em}Small & 48,306 (1.3\%)\\
\hspace{1em}Medium & 59,979 (1.6\%)\\
\hspace{1em}Large & 3,484,092 (95\%)\\
\addlinespace
Public transit availability in the neighbourhood & 3,537,339 (96\%)\\
Montly or annual public transit pass & 1,083,351 (29\%)\\
\bottomrule
\multicolumn{2}{l}{\rule{0pt}{1em}\textsuperscript{1} Mean (SD); n (\%)}\\
\end{longtable}
\endgroup{}

\subsubsection{Controlling for
endogeneity}\label{controlling-for-endogeneity}

In social science, it is well understood that \emph{``correlation does
not imply causality''} \citep{cunningham2021}. It is important to
evaluate the actual impact of a predictor variable (the cause) on the
outcome (the effect) when modeling the causal inference. Including
causal inference in our modeling approach means we must account for the
possible (and indeed likely) presence of endogeneity, which occurs when
an independent variable correlates with the regression error term
\citep{legallo2013} and compromises the consistency of the estimators
\citep{antonakis2010}. In this study, we propose to address endogeneity
by applying an instrumental variable (IV) technique. An instrumental
variable \(Z\) must satisfy the following criteria
(\citep{cunningham2021, luz2022}):

\begin{enumerate}
\def\labelenumi{\arabic{enumi}.}
\tightlist
\item
  \(Z\) must be highly correlated with and have a causal effect on the
  endogenous variable \(D\).
\item
  \(Z\) must affect the dependent variable \(Y\) only through \(D\),
  with no direct effect of \(Z\) on \(Y\).
\item
  \(Z\) must not be correlated with the regression residuals of \(D\) on
  \(Y\).
\end{enumerate}

In simpler terms, the instrument must be correlated with the outcome
only through accessibility. Examples of instrumental variables used in
studies analyzing the relationship between employment outcomes and job
accessibility include Euclidean distance from the spatial unit centroid
to the nearest major road or employment subcenter
(\citep{delmelle2021, jin2018}), share of no-vehicle households
(\citep{hu2017}), transit accessibility and car unavailability
(\citep{johnson2017}), distance to rivers ({[}\citet{duarte2023};
Greg´orio Luz 2022{]}), and population density
({[}\citet{bastiaanssen2021}; Bastiassen 2026{]}).

However, two challenges arise in the search for relevant instruments: i)
the likelihood that the instruments and error terms are correlated
increases as the correlation between the endogenous variable and the
instruments becomes stronger, which can render the instruments invalid;
and ii) instruments that are only weakly correlated with the endogenous
variable may be ineffective and perform poorly. In response to these
challenges, Le Gallo and Páez (\citeyearpar{legallo2013}) developed an
approach that relies on synthetic instrumental variables when the model
involves spatial data. The researchers demonstrated that synthetic
variables can be generated using spatial filtering, a technique that
removes spatial residual autocorrelation to improve regression analysis.
They propose the use of an eigenvector-based technique, in which
transforms a projection of the contiguity matrix into latent map
patterns through eigenvector structural analysis. These latent patterns,
included in the spatial structure of the system, are then used to
generate synthetic variables, which serve as instruments in IV
estimation. Since they are derived from the spatial structure of the
system, these synthetic instruments are exogenous, thereby fulfilling
the core requirement of IV estimation.

Following the methodological approach of Gallo and Páez
\citeyearpar{legallo2013} to create spatial filters for the
accessibility variables:

\begin{enumerate}
\def\labelenumi{\arabic{enumi}.}
\tightlist
\item
  We built a queen-contiguity matrix for the census tract to represent
  spatial weights, expressing the existence of a neighbour relation as a
  binary relationship (1 for neighbours and 0 otherwise).
\item
  Then, we computed all eigenvectors associated with the contiguity
  matrix. These eigenvectors are orthogonal to one another and represent
  a portion of the variance in the spatial contiguity matrix.
\item
  After this, we initialized a variable \(S\) with value \(0\) and
  linearly regressed each accessibility indicator (the independent
  variable) based on each eigenvector and \(S\). If the eigenvector
  coefficient was significant (adopting a p-value \(\leq\) 0.05), we
  multiplied the eigenvector by its coefficient and added to value of
  \(S\). After this process, \(S\) represents a spatial filter for the
  accessibility indicator.
\end{enumerate}

All spatial filters presented a Pearson correlation of above 0.88 with
their respective accessibility measures, and, generally, absolute
correlations lower than 0.1 with the original perception dependent
variables (before transforming them into categorical variables, Appendix
A). Since our dependent variables are not linear, the popular technique
of replacing endogenous variables by their first-stage predictors
(two-stage least square estimation) is not applicable (). Therefore, we
applied the two-stage residual inclusion (2SRI) approach. In the
first-stage, we estimated the accessibility measures using a linear
regression of all control variables plus spatial filters (our
instrumental variables). In the second-stage, the dependent variables
were estimated using the accessibility measures, the control variables,
and the residuals from the first-stage regressions. In both stages, the
respondent weights were considered to perform the regressions.

\subsection{Equity Standards}\label{equity-standards-1}

We selected the model that best explains the effect of calculated
accessibility on the perceptions of accessibility to obtain equity
standards. We measured changes in the probabilities across perception
categories (from most negative to most positive) by varying the levels
of calculated accessibility, while controlling for the other explanatory
variables. Controlling for the explanatory variables means holding their
values constant, while varying the values of the treatment. Based on
this, we decided to evaluate the effects of calculated accessibilities
variables in three distinct population profiles:

\begin{itemize}
\tightlist
\item
  An equity-deserving population group, with individual-level control
  variables set to the values corresponding to populations facing the
  greatest transportation challenges.
\item
  The most common population, reflecting modal characteristics of the
  study area population and with control variables set to their most
  common values; and,
\item
  The advantaged group, defined by selecting for the most privileged
  population segment in terms of their transport experience.
\end{itemize}

For each population profile, we increased the value of each
accessibility measure by 0.01 jobs per person, while holding constant
other accessibility measures (representing alternative transportation
modes) at their weighted median values. We also adopted the weighted
median value for the population density and a constant zero value for
the residuals obtained from the first-stage regressions. Table X
presents the control variable values applied for each population group.

Based on the best model and for each population group, we identified the
minimum value of calculated job accessibility per person in which the
probability of the positive and neutral perceptions combined was higher
than the probability of being in negative perceptions, naming this
minimum value as a conservative threshold. Next, we defined a
progressive threshold as the minimum accessibility value that resulted
in the probability of positive perceptions (not considering neutral
perceptions) being higher than of negative perceptions. We then
identified an intermediate threshold, set as the midpoint between the
conservative and progressive thresholds.

Then, we selected as the equity standards the accessibility thresholds
obtained based on the equity-deserving population. The rationale behind
this equity standard definition is to identify the minimum accessibility
level that gives the most disadvantaged population a higher probability
of having at least a neutral perception (conservative standard), a
neutral to positive perception (intermediate standard), or a positive
perception (progressive standard).

\subsection{Measuring transportation
equity}\label{measuring-transportation-equity}

We calculated the extent and severity of accessibility gap using
Foster-Greer-Thorbecke (FGT) measures, according to the Equation
\ref{eq:FGT}. FGT measures evaluates the number of people below a
standard and the size of the accessibility shortfall relative to it.

\begin{equation}
\label{eq:FGT}
FGT_\alpha = \frac{1}{N}\sum_{H}^{i=1}\left(\frac{z-y_i}{z}  \right)^{\alpha}
\end{equation}

Where \(N\) is 0the size of the total population; \(z\) is the
accessibility standard; \(H\) is the number of individuals below the
\(z\); \(y_i\) is the accessibility level of each individual i; and
\(\alpha\): determines how sensitive the index is to change in the
accessibility gap.

Changing the values of \(\alpha\) enables the calculation of different
FGT indices. When \(\alpha = 0\), \(FGT_0\) measures the extent of
undersupply as a simple headcount, indicating the proportion of people
below the standard. When \(\alpha = 1\), \(FGT_1\) measures the severity
of the undersupply as the average percent distance between the standard
and the accessibility levels of the individuals below it. When
\(\alpha = 2\), \(FGT_2\) measures both the extent and severity of
undersupply by calculating the number of people below the standard,
weighted by the size of the accessibility shortfall relative to the
standard.

\subsection{Areas for priority
intervention}\label{areas-for-priority-intervention}

For each equity standard, we identified CTs with accessibility shortfall
(cases in which the CT accessibility was lower than the equity
threshold). Next, we calculated the Accessibility Fairness Index (AFI)
(Martens, 2017) to measure the severity of the insufficiency, by
multiplying the population size below the standard by the accessibility
gap. In practice, the AFI is a simple way to highlight areas with large
groups of people that fall the furthest behind the set standards (van
der Veen et al., 2020). After this, we classified CTs with accessibility
shortfall into levels of priority for policy interventions: Very-high
(AFI above the third quartile, top 25\%), High (AFI between the median
and third quartile, 25-50\% from the top), Moderate (AFI between the
first quartile and the median, 50-75\% from top), and Low priority (AFI
below the first quartile, bottom 25\%).

\section{Results and discussions}\label{results-and-discussions}

\subsection{The effect of calculated accessibility on perceived
accessibility}\label{the-effect-of-calculated-accessibility-on-perceived-accessibility}

\begin{landscape}\begin{table}[!h]
\centering
\caption{\label{tab:unnamed-chunk-4}\label{tab:models}Models.}
\centering
\resizebox{\ifdim\width>\linewidth\linewidth\else\width\fi}{!}{
\begin{tabular}[t]{lcccccccc}
\toprule
\multicolumn{1}{c}{ } & \multicolumn{2}{c}{\textbf{Few jobs}} & \multicolumn{2}{c}{\textbf{Easy to reach}} & \multicolumn{2}{c}{\textbf{Travel too long}} & \multicolumn{2}{c}{\textbf{Decline employment}} \\
\cmidrule(l{3pt}r{3pt}){2-3} \cmidrule(l{3pt}r{3pt}){4-5} \cmidrule(l{3pt}r{3pt}){6-7} \cmidrule(l{3pt}r{3pt}){8-9}
\textbf{Variable} & \textbf{OR} & \textbf{SE} & \textbf{OR} & \textbf{SE} & \textbf{OR} & \textbf{SE} & \textbf{OR} & \textbf{SE}\\
\midrule
Active transportation job accessibility per user & 0.86* & 0.064 & 0.85 & 0.142 & 0.95 & 0.157 & 0.84* & 0.081\\
Transit job accessibility per user & 1.25* & 0.112 & 0.92 & 0.132 & 0.93 & 0.145 & 0.88 & 0.184\\
Car job accessibility per user & 0.93 & 0.097 & 1.08 & 0.119 & 1.11 & 0.111 & 1.18 & 0.139\\
Number of Cars & 0.98 & 0.121 & 1.06 & 0.148 & 1.04 & 0.158 & 1.11 & 0.136\\
Number of Bikes & 1.08 & 0.067 & 0.95 & 0.093 & 0.95 & 0.074 & 1.02 & 0.065\\
\addlinespace
Preferred mode match most frequent mode & 1.14 & 0.168 & 1.72* & 0.243 & 1.68* & 0.225 & 1.30 & 0.246\\
Montly or annual public transit pass & 0.68* & 0.174 & 1.35 & 0.220 & 0.86 & 0.251 & 0.52* & 0.256\\
Public transit availability in the neighbourhood & 0.75 & 0.245 & 0.97 & 0.361 & 1.11 & 0.301 & 1.06 & 0.226\\
Indigenous or racialized person & 0.60** & 0.164 & 0.72 & 0.218 & 0.71 & 0.221 & 0.57** & 0.210\\
Temporary visa or refugee & 1.61 & 0.257 & 1.00 & 0.400 & 0.71 & 0.359 & 0.71 & 0.414\\
\addlinespace
Person with disability & 0.69 & 0.209 & 0.89 & 0.275 & 1.02 & 0.312 & 0.80 & 0.229\\
Woman, transgender or non-binary person & 1.05 & 0.149 & 0.84 & 0.188 & 1.45 & 0.203 & 1.03 & 0.194\\
Age & 1.01 & 0.006 & 1.00 & 0.008 & 1.01 & 0.008 & 1.02* & 0.007\\
Household size & 1.06 & 0.052 & 0.96 & 0.078 & 1.09 & 0.087 & 0.94 & 0.067\\
Household income &  &  &  &  &  &  & ** & \\
\addlinespace
\hspace{1em}Less than \$40,000 & — & — & — & — & — & — & — & —\\
\hspace{1em}\$40,000 to \$80,000 & 0.80 & 0.188 & 1.00 & 0.257 & 0.78 & 0.247 & 1.48 & 0.269\\
\hspace{1em}\$80,000 to \$120,000 & 1.23 & 0.220 & 1.07 & 0.257 & 0.62 & 0.281 & 1.96 & 0.248\\
\hspace{1em}More than \$120,000 & 1.29 & 0.227 & 1.37 & 0.312 & 0.50 & 0.305 & 3.08 & 0.312\\
Population density (people/km2) & 1.18* & 0.075 & 1.13 & 0.116 & 1.28* & 0.103 & 1.10 & 0.100\\
\addlinespace
Population centre class &  &  &  &  &  &  & *** & \\
\hspace{1em}Rural & — & — & — & — & — & — & — & —\\
\hspace{1em}Small & 1.34 & 0.486 & 0.68 & 0.856 & 1.59 & 0.815 & 0.89 & 0.478\\
\hspace{1em}Medium & 2.35 & 0.469 & 1.41 & 0.583 & 3.16 & 0.571 & 5.06 & 0.452\\
\hspace{1em}Large & 1.24 & 0.322 & 1.33 & 0.438 & 1.34 & 0.475 & 1.16 & 0.330\\
\addlinespace
Residuals: Active acessibility & 1.09 & 0.078 & 1.01 & 0.128 & 1.03 & 0.137 & 0.97 & 0.085\\
Residuals: Transit acessibility & 0.85 & 0.092 & 1.49** & 0.130 & 1.32* & 0.118 & 0.85 & 0.146\\
Residuals: Car acessibility & 1.13 & 0.068 & 0.91 & 0.127 & 1.06 & 0.099 & 0.98 & 0.094\\
Active mode: accessibility × user & 1.17* & 0.072 & 1.10 & 0.342 & 1.32 & 0.359 & 1.20* & 0.088\\
Transit mode: accessibility × user & 0.95 & 0.133 & 1.03 & 0.169 & 0.83 & 0.183 & 1.17 & 0.205\\
\bottomrule
\multicolumn{9}{l}{\rule{0pt}{1em}Abbreviations: CI = Confidence Interval, OR = Odds Ratio, SE = Standard Error}\\
\multicolumn{9}{l}{\textsuperscript{} *p<0.05; **p<0.01; ***p<0.001}\\
\end{tabular}}
\end{table}
\end{landscape}

\subsection{Standards}\label{standards}

\subsection{Areas for priority
intervention}\label{areas-for-priority-intervention-1}

\section{Limitations}\label{limitations}

\section{Summary and conclusions}\label{summary-and-conclusions}

\renewcommand\refname{References}
\bibliography{mybibfile.bib}


\end{document}
